% New analytic increment. No preceding increment is modified.
\documentclass[11pt]{article}
\usepackage[T1]{fontenc}
\usepackage{lmodern,amsmath,amssymb,amsthm,mathtools,microtype,booktabs,xurl}
\usepackage[margin=.9in]{geometry}
\usepackage[hidelinks]{hyperref}
\setlength{\emergencystretch}{3em}
\allowdisplaybreaks
\numberwithin{equation}{section}
\newtheorem{theorem}{Theorem}[section]
\newtheorem{lemma}[theorem]{Lemma}
\newtheorem{proposition}[theorem]{Proposition}
\newtheorem{corollary}[theorem]{Corollary}
\theoremstyle{remark}\newtheorem{remark}[theorem]{Remark}
\newcommand{\E}{\mathbb E}
\newcommand{\dd}{\,d}
\newcommand{\one}{\mathbf1}
\newcommand{\Sp}{\mathcal S_\pi}
\newcommand{\Ss}{\mathcal S_{3\pi/2}}
\newcommand{\Kh}{\mathcal K}
\title{The $\pi$-strip budget and middle-zone freezing for both Q3 strips}
\author{Analytical increment for independent review}
\date{Version 1; October 22, 2026}
\begin{document}
\maketitle
\begin{abstract}
On the existing S2CP/CFA stopped interval, without CFA's strip input
or its passage conclusion, both required $\pi$-strip mass budgets
vanish. We prove the refined middle-zone freezing estimate for both
strips. For the inner $\pi$ zone, scalar target-only order bounds its
mass at the moving overlap by the regularized boundary trajectory.
A coarse, explicitly paid strong-rate bound then suffices until
half mass. Combined with the capped weak gain, its remaining uniform
power is $q^{411/800}$ in mass divided by $q$. This is weaker than
the proposed $q^3(\log\ell)^{3/2}$ mass rate, but closes both requested
budgets with no additional hypothesis. No new labels are put into
$G_+$, and the $3\pi/2$ inner zone is left for its stated consumer.
\end{abstract}

\section{Setting and the resulting budget}
All new labels have prefix \texttt{an03pisbv1:}.
Use the original initialized balanced flow
\[
 U(0,b)=W(0,b)=q^5(\cos b,\sin b),\quad
 U=\sqrt m\,a_\theta,\quad W=\sqrt m\,a_\psi,\quad
 \dd\lambda_0=\dd b/(2\pi),
\]
\[
 P=(P_1+P_2)/2,\quad P_1=N(e_1,q^2I_2),\quad
 P_2=N(\rho e_2,q^2I_2),\quad
 \rho\in[13/20,7/10],\quad \lambda=\rho^2,\quad Q=\lambda+q^2 .
\]
Let $f(X)=\int W(U\cdot X)_+\dd\lambda_0$, $e_i=X_i-f_i$,
$R_{ij}(\theta)=\E_P[e_iX_j\one_{\{a_\theta\cdot X>0\}}]$.
The exact equations used below are
\begin{equation}
 \theta'=a_\psi^TR_\theta a_\theta^\perp,\qquad
 \psi'=(a_\psi^\perp)^TR_\theta a_\theta,\qquad
 (\log m)'=2a_\psi^TR_\theta a_\theta .
 \label{an03pisbv1:exact}
\end{equation}
Set
\begin{align}
 \ell&=\log(1/q),\quad \omega=(1-\lambda)/2,\quad
 c_m=2/\omega+1/2,\quad u=10-c_m\ge169/102,\nonumber\\
 t_m&=c_m\ell,\qquad
 t_0=t_h=10\ell-\log A_q+O(q^2\ell^2+q^u),\quad
 0<c_A\le A_q\le C_A,\nonumber\\
 k&=10(1-\lambda),\qquad \beta=\lambda/Q,\qquad
 I(t)=Q\int_{t_0}^t(1-c),\quad
 c=Q^{-1}\E_{P_2}[f_2X_2].                       \label{an03pisbv1:scales}
\end{align}
We work on any stopped interval $[t_0,t_*]$ of S2CP with CFA's
retained non-strip inputs. In particular $0\le c\le c_b<1$,
$t_*-t_0\le L_T\ell$, $\int r<B$,
$r=\sup_\theta|R_{11}(\theta)|$, the cross residual is at most
$M_{\rm res}q$, and the primary energy is below $E_g$.
Its instantaneous core hypotheses and the joint secondary guard or
proved improvement remain in force. We do \emph{not} assume CFA's
strip supremum/integral input or apply its conditional passage.
The existing outside supremum, wherever used by S2CP's gate lemma,
is still a stopped input, not a new conclusion about the other strip.

The proved inputs needed from S1 and S2 are recorded explicitly:
\begin{align}
 \sup_{[0,t_0]}(\|e_1\|_{P_2}+\|e_2\|_{P_1})
       &\le C_{\rm pre}q,                         \label{an03pisbv1:cross}\\
 \int_{t_0}^{t_*}X(t)\dd t
       &\le G+o(1),\quad
 X=\sup_\theta(|R_{22}|-I'/2)_+,\quad
 G=\frac{5eE_g}{9(1-c_b)},                         \label{an03pisbv1:gate}\\
 e^{I(t)}
       &\le C_H E_g A_q^\lambda q^{-k}e^{D_{\rm clk}},
 \quad 0\le D_{\rm clk}\le C(q^{1/8}+q^2\ell).
                                                        \label{an03pisbv1:clock}
\end{align}
Equation \eqref{an03pisbv1:cross} is Q3SB's proved pre-entry result.
Equations \eqref{an03pisbv1:gate}--\eqref{an03pisbv1:clock} are S2CP's
gate lemma and its stopped clock, using the original Q2 clock seed
$H(t_0)\asymp A_q^{-\lambda}q^k$. They are conclusions under the
current guards, not hypotheses of a later passage theorem.
Loss monotonicity supplies the remaining diagonal residual norms
on $[0,t_*]$. MOMC gives $t_*\le C_{\rm dur}\ell$.

Fix $0<\varepsilon<1/2$ before $q$ and define
\begin{equation}
 \begin{gathered}
 k_G=169/3200,\quad \Lambda_q=\sqrt{3\ell/k_G},\quad
 w_q=\arcsin(q\Lambda_q),\quad
 s_q=\sqrt{(1+\varepsilon)\log\ell/k_G},\\
 \Sp=[\pi,\pi+2w_q),\qquad
 \Ss=(3\pi/2-2w_q,3\pi/2].
 \end{gathered}                                   \label{an03pisbv1:zones}
\end{equation}
In either strip let $a$ be its initial nonnegative physical offset
from its axis: $a=b-\pi$ or $a=3\pi/2-b$.
The inner zone is $a<qs_q$, the middle zone is
$qs_q\le a\le(1+\varepsilon)w_q$, and the outer zone is
$a>(1+\varepsilon)w_q$.

\begin{theorem}[The $\pi$ strip is paid; both other zones freeze]
\label{an03pisbv1:main}
For the stated original flow and stopped interval, the middle and
outer zones of each strip have
\begin{equation}
 \sup_{[t_0,t_*]}\frac{\mu_{\rm mid\cup out}}q
       \le Cq^{10}\sqrt\ell,\qquad
 \int_{t_0}^{t_*}\frac{\mu_{\rm mid\cup out}}q
       \le CL_Tq^{10}\ell^{3/2}=o(1).
 \label{an03pisbv1:frozenbudget}
\end{equation}
For the whole $\pi$ strip put
\begin{align}
 \delta_1(\lambda)&=c_m-10(1-\lambda),\nonumber\\
 \delta_2(\lambda)&=(19\lambda-7)/2,\nonumber\\
 \delta_3(\lambda)&=4/\lambda+10\lambda-12,\qquad
 \delta_*=\frac{411}{800}.                         \label{an03pisbv1:powers}
\end{align}
All three $\delta_j\ge\delta_*>0$. There is an explicit vanishing
error bound
\[
 \Xi_q=D_{\rm clk}
       +C\{q\ell+q^2\ell^2+q^u+(1+M_{\rm res})qL_T\ell\}=o(1)
\]
such that
\begin{equation}
 \sup_{[t_0,t_*]}\frac{\mu_{\Sp}}q
 \le Cq^{10}\sqrt\ell+
     C e^{2B+\Xi_q}E_g A_q^{\lambda-1}s_q
                     \sum_{j=1}^3q^{\delta_j(\lambda)}
 =o(1).
 \label{an03pisbv1:pibudget}
\end{equation}
Multiplying its right side by $L_T\ell$ bounds
$\int_{t_0}^{t_*}\mu_{\Sp}/q$, which is also $o(1)$.
In particular the rate obtained uniformly on the ratio box is
\[
 \mu_{\Sp}
   =O\!\left(q^{1211/800}\sqrt{\log\ell}
                   +q^{11}\sqrt\ell\right)
\]
with the fixed factors in \eqref{an03pisbv1:pibudget} retained.
The sharper proposed $O(q^3(\log\ell)^{3/2})$ rate is not claimed.
No $G_+$ enlargement or new entry contract is required.
\end{theorem}

\section{Outer and middle freezing}
\begin{lemma}[The full outer buffer]
\label{an03pisbv1:outer}
There is a fixed $C_{\rm buf}$ such that every initial label in
$(\pi+w_q,3\pi/2-w_q)$ at distance at least
$C_{\rm buf}q^3\ell$ from both ends stays in that interval on
$[0,t_*]$ and satisfies
\[
 |\theta(t,b)-b|+|\psi(t,b)-b|\le Cq^3\ell,\qquad
 m(t,b)=q^{10}e^{O(q^3\ell)}.
\]
This applies to both outer zones in \eqref{an03pisbv1:zones}.
\end{lemma}
\begin{proof}
Until the first exit, both input projected means are at most
$-q\Lambda_q$ in the coordinate convention of GESG interiorspeed.
That lemma, valid from initialization by loss alone, bounds the
two angular speeds and the logarithmic radial speed by
$Ce^{-k_G\Lambda_q^2}=Cq^3$.
Choose $C_{\rm buf}>2CC_{\rm dur}$. The available travel is smaller
than the initial buffer, so no first exit exists.
The outer-zone buffers are at least $\varepsilon w_q$ at the
nearby end and a fixed positive constant at the other end, and
$\varepsilon w_q\gg q^3\ell$.
\end{proof}

\begin{lemma}[Every selected matrix entry retains its small factor]
\label{an03pisbv1:matrix}
Let $T(s)=e^{-k_Gs^2}$ and
$C_*=C_{\rm model}(1+C_{\rm pre}+M_{\rm res})$.
For $0\le s_\theta,s_\psi\le3\Lambda_q$, the following bounds
hold throughout $[0,t_*]$.
Near $\pi$, use $\theta=\pi+qs_\theta$, $\psi=\pi+qs_\psi$;
near $3\pi/2$, use $\theta=3\pi/2-qs_\theta$,
$\psi=3\pi/2-qs_\psi$. Entrywise,
\begin{equation}
 \begin{array}{c|c}
 \text{near }\pi&\text{near }3\pi/2\\[2pt]\hline
 |R_\theta|\le C_*T(s_\theta)
       \begin{pmatrix}q^2&q\\q&1\end{pmatrix}
 &
 |R_\theta|\le C_*T(s_\theta)
       \begin{pmatrix}1&q\\q&q^2\end{pmatrix}.
 \end{array}                                      \label{an03pisbv1:table}
\end{equation}
Consequently
\begin{align}
 |\theta'|+|\psi'|
 &\le C_*q(1+s_\theta+s_\psi)T(s_\theta),
                                                        \label{an03pisbv1:angular}\\
 |(\log m)'|
 &\le C_*q^2(1+s_\theta)(1+s_\psi)T(s_\theta).
                                                        \label{an03pisbv1:radial}
\end{align}
The $C_*$ in each display may be enlarged by a model-only factor.
\end{lemma}
\begin{proof}
For any cluster and selector $A$, Cauchy--Schwarz and Holder give
\[
 |\E[e_iX_j\one_A]|
 \le \|e_i\|_2(\E X_j^4)^{1/4}\mathbb P(A)^{1/4}.
\]
Near $\pi$, the $P_1$ gate is exponentially closed in $q^{-2}$.
On $P_2$, the fourth-moment norms of $X_1,X_2$ are respectively
$O(q),O(1)$, while the residual norms of $e_1,e_2$ are respectively
$O((C_{\rm pre}+M_{\rm res})q),O(1)$.
This proves the four powers in the left matrix.
Near $3\pi/2$, the $P_2$ gate is exponentially closed.
On $P_1$, the fourth-moment norms are $O(1),O(q)$ and the residual
norms are $O(1),O((C_{\rm pre}+M_{\rm res})q)$.
This proves all four powers in the right matrix.
In particular no smallness of $R_{11}$ is assumed there.

The nontrivial selected gate on $P_2$ has probability at most
$\exp[-\rho^2\sin^2(qs_\theta)/(2q^2)]$, and the corresponding
$P_1$ bound has $\rho^2$ replaced by one. For
$s_\theta\le3\Lambda_q$,
\[
 \frac{\rho^2\sin^2(qs_\theta)}{8q^2}
 \ge k_Gs_\theta^2-Cq^2s_\theta^4 .
\]
The error is $O(q^2\ell^2)=o(1)$, so a fixed factor suffices to
retain the exact coefficient $k_G$. The other cluster's
$e^{-c/q^2}$ contribution is at most $Cq^2T(s_\theta)$:
on this domain $T(s_\theta)\ge q^{27}$.
Thus every entry, including the noise-coordinate entries, retains
the same gate tail.

For the angular equations insert \eqref{an03pisbv1:table} in
\eqref{an03pisbv1:exact}. Near $\pi$, the two sines have magnitudes
at most $qs_\theta,qs_\psi$; near $3\pi/2$ the same holds for the
cosines. This proves \eqref{an03pisbv1:angular}.
For the radial equation the large diagonal is multiplied by both
small coordinates, each cross entry by one, and the remaining
diagonal is already $O(q^2)$. This proves \eqref{an03pisbv1:radial}.
\end{proof}

\begin{proposition}[Uniform freezing in the middle zone]
\label{an03pisbv1:middle}
For a middle-zone initial scaled offset $s_0=a/q$, uniformly on
$[0,t_*]$,
\begin{align}
 |s_\theta(t)-s_0|+|s_\psi(t)-s_0|
 &\le C\ell(1+s_0)e^{-k_Gs_0^2}
 \le C(1+s_q)\ell^{-\varepsilon}=o(1),             \label{an03pisbv1:travel}\\
 \left|\log\frac{m(t,b)}{q^{10}}\right|
 &\le Cq^2\ell(1+s_0)^2e^{-k_Gs_0^2}
 \le Cq^2(1+s_q)^2\ell^{-\varepsilon}=o(1).
                                                        \label{an03pisbv1:middleradial}
\end{align}
Here $s_0$ is an offset, not RCET's initial sine.
\end{proposition}
\begin{proof}
For a sufficiently large fixed $A$, set
$\Delta(s_0)=A\ell(1+s_0)e^{-k_Gs_0^2}$.
Uniformly for $s_0\ge s_q$ in the middle zone,
\[
 \Delta(s_0)=o(1),\qquad
 s_0\Delta(s_0)\le C(1+s_q)^2\ell^{-\varepsilon}=o(1).
\]
These bounds follow because the polynomial times Gaussian is
decreasing on this range for small $q$.
Stop when the sum of the two scaled displacements first equals
$\Delta(s_0)$. Before that time, the offsets are positive and below
$3\Lambda_q$, and
\[
 T(s_\theta)\le T(s_0)e^{2k_Gs_0\Delta(s_0)}
                      \le 2T(s_0)
\]
for small $q$. The scaled angular bound integrated over
$C_{\rm dur}\ell$ is at most
$C C_{\rm dur}\ell(1+s_0)T(s_0)<\Delta(s_0)/2$
after choosing $A$. This contradicts the first equality.
Integrating the radial bound on the same region gives
\eqref{an03pisbv1:middleradial}.
A label-dependent $\Delta(s_0)$ avoids discarding a gate-tail
constant when $s_0$ is as large as $\Lambda_q$.
\end{proof}

The union of the two frozen zones in either strip has initial width
at most $2w_q=O(q\sqrt\ell)$. Lemma~\ref{an03pisbv1:outer} and
Proposition~\ref{an03pisbv1:middle}, with the fixed label measure
$\dd b/(2\pi)$, prove \eqref{an03pisbv1:frozenbudget}.
This proves tasks (1)--(2) without any new entry profile.

\section{The inner \texorpdfstring{$\pi$}{pi} zone at the moving overlap}
Hats denote only the exact aligned target-only flow with the same
initial label and norm. Let $\Kh(z)=\phi(z)+z\Phi(z)$.
Its angular and radial equations are
\[
 \widehat\theta'=v_q(\widehat\theta),\qquad
 v_q(\vartheta)=\frac q2[-\sin\vartheta\,\Kh(\cos\vartheta/q)
       +\rho\cos\vartheta\,\Kh(\rho\sin\vartheta/q)],
\]
\[
 (\log\widehat m)'=2g_q(\widehat\theta),\qquad
 g_q(\vartheta)=\E_P[(a_\vartheta\cdot X)_+^2],\qquad
 g_q'=2v_q .
\]

\begin{lemma}[Boundary domination in the target flow]
\label{an03pisbv1:domination}
For $\pi\le b\le\pi+qs_q$ and every target-only time $t\ge0$,
\begin{equation}
 \widehat\theta(t,b)\ge\widehat\theta(t,\pi),\qquad
 \widehat m(t,b)\le\widehat m(t,\pi).
 \label{an03pisbv1:targetorder}
\end{equation}
For the original flow at $t_m$, the proved all-label early comparison
therefore gives
\begin{equation}
 m(t_m,b)\le e^{Cq^u}\widehat m(t_m,\pi),\qquad
 \theta(t_m,b),\psi(t_m,b)
       \ge\theta(t_m,\pi)-Cq^u .
 \label{an03pisbv1:matching}
\end{equation}
No order or cooperativity of the finite-noise original field is
assumed.
\end{lemma}
\begin{proof}
PROFILE gives its attracting root $qa_q$ and $v_q<0$ above that
root on a fixed strong chart. The sign continues through Q1:
using $\Kh(z)=z+\Kh(-z)$,
\[
 v_q=-\omega\sin\vartheta\cos\vartheta
 +\frac q2[-\sin\vartheta\,\Kh(-\cos\vartheta/q)
          +\rho\cos\vartheta\,\Kh(-\rho\sin\vartheta/q)]<0
\]
on $qL\le\vartheta<\pi/2$ for a sufficiently large fixed $L$;
use $\Kh(-z)\le\Kh(0)$ and $\sin\vartheta\ge2qL/\pi$.
At $\pi/2$ the negative first term in the exact formula remains.
On $[\pi/2,\pi]$ both terms in the exact formula are nonpositive,
and their sum is strictly negative.

For $\vartheta=\pi+qs$, $0\le s\le s_q$, the only positive
contribution has size at most $Cq^2s_q e^{-c/q^2}$.
The negative contribution has magnitude at least
\[
 cq\,\Kh(-\rho\sin(qs)/q)
       \ge \frac{cq\,e^{-\rho^2s_q^2/2}}{(1+s_q)^2}.
\]
The last bound follows by integrating the Gaussian definition
$\Kh(-z)=\int_z^\infty(v-z)\phi(v)\dd v$
over an interval of length $(1+z)^{-1}$ above $z$.
Its decay is only a fixed power of $\ell$ times a logarithmic
factor, so it dominates the exponentially closed first gate.
Thus $v_q<0$ throughout $(qa_q,\pi+qs_q]$.

Scalar uniqueness preserves the angular order, and both
characteristics stay in this interval. Since $g_q'=2v_q<0$,
their logarithmic radial rates are ordered as well. Equal initial
masses prove \eqref{an03pisbv1:targetorder}.
Q2EC's equation early is explicitly valid for every initial label
at $t_m$; apply it to $b$ and to $\pi$ to obtain
\eqref{an03pisbv1:matching}. Its angle error is retained; no exact
original-flow angular order is inferred.
\end{proof}

\begin{lemma}[A regularized upper bound at the boundary label]
\label{an03pisbv1:boundary}
Uniformly on the ratio box,
\begin{equation}
 \widehat m(t_m,\pi)
 \le C e^{Cq\ell}
   \left(q^{10}+q^{12-\lambda c_m}
                  +q^{8-c_m+4/\lambda}\right).
 \label{an03pisbv1:boundarymass}
\end{equation}
\end{lemma}
\begin{proof}
PROFILE's exact target coordinate equations are
\[
 \widehat U_j'
  =\frac{\mu_jq|\widehat U|}{2}
       \Kh(\mu_j\widehat U_j/(q|\widehat U|))+d_q\widehat U_j,
 \quad (\mu_1,\mu_2)=(1,\rho),\quad 0\le d_q\le q^2.
\]
For $b=\pi$, initially $\widehat U_1=-q^5$, $\widehat U_2=0$.
Let $T_\times$ be the first $\widehat U_1=0$.
Before this crossing, $\widehat U_2\ge0$ and
$|\widehat U_1|\le q^5e^{q^2t}$.
Using $\Kh(z)\le z+\Kh(0)$ for $z\ge0$ gives
\[
 \widehat U_2'
 \le(\lambda/2+Cq+q^2)\widehat U_2+Cq^6e^{q^2t},
 \qquad
 \widehat U_2(t)\le Cq^6e^{\lambda t/2+Cqt+q^2t}.
\]
Thus before crossing
$\widehat m(t,\pi)\le C(q^{10}+q^{12}e^{\lambda t})e^{Cqt}$
on logarithmic times.

Q2E's endpoint-uniform crossing formula, also quoted in RCET,
retains $(s_0+q)$ and $(c_0+q)$ and gives at this boundary
\[
 T_\times=\frac4\lambda\ell+O(1),\qquad
 \widehat m(T_\times,\pi)\asymp q^8 .
\]
The $O(1)$ is uniform on the ratio box.
After crossing, the target covariance bound gives
$(\log\widehat m)'\le1+2q^2$.
If $T_\times\le t_m$, its contribution is therefore at most
$Cq^{8-c_m+4/\lambda}e^{Cq^2\ell}$.
If $T_\times>t_m$, the pre-crossing bound applies instead.
The sum of the three nonnegative terms in
\eqref{an03pisbv1:boundarymass} covers both cases.
In particular the Gaussian seed $q^6$ in $\widehat U_2$, and
hence the $q^{12}e^{\lambda t}$ term, has not been dropped.
\end{proof}

\section{Paying the remaining pre-entry growth without a new Q2 contract}
\begin{lemma}[A coarse all-label bridge to half mass]
\label{an03pisbv1:bridge}
There is a model/S1 constant $C$ such that for every inner
$\pi$ label
\begin{equation}
 m(t_0,b)\le C A_q^{-1}
     e^{C(q\ell+q^2\ell^2+q^u)}
       \left(q^{c_m}
          +q^{2+(1-\lambda)c_m}
          +q^{4/\lambda-2}\right).
 \label{an03pisbv1:entrymass}
\end{equation}
This is only a mass bound. It asserts no new Q2EC ratio, sign,
antisymmetry, or secondary profile contract for these labels.
\end{lemma}
\begin{proof}
Let $S(t)=\int m(t,b)\dd\lambda_0(b)$. For every receiver write
$R_\theta=A_\theta-C_\theta$, with
\[
 A_\theta=\E_P[XX^T\one_{\{a_\theta\cdot X>0\}}],\qquad
 C_\theta=\E_P[fX^T\one_{\{a_\theta\cdot X>0\}}].
\]
The target matrix is positive semidefinite and
$\|A_\theta\|\le\|\E_PXX^T\|=(1+2q^2)/2$.
Also $|f(X)|\le S|X|$, so for $q\le1/2$,
$2\|C_\theta\|\le2S\E_P|X|^2\le4S$.
Consequently, with no angular restriction,
\begin{equation}
 (\log m)'\le1+2q^2+4S(t).                         \label{an03pisbv1:coarserate}
\end{equation}
On $[t_m,t_0]$, MOMC, as explicitly recorded in Q2EC's
equation inputs, supplies
\[
 S=M+\mu,\quad \int M\le3/2,\quad
 \mu\le Cq^2\ell^2M.
\]
Thus $\int S\le3/2+Cq^2\ell^2$ and
\[
 m(t_0,b)\le
 e^{6+O(q^2\ell^2)}e^{(1+2q^2)(t_0-t_m)}m(t_m,b).
\]
Insert \eqref{an03pisbv1:matching} and
\eqref{an03pisbv1:boundarymass}, and use
$e^{t_0-t_m}=A_q^{-1}q^{c_m-10}
 e^{O(q^2\ell^2+q^u)}$. This proves
\eqref{an03pisbv1:entrymass}.
\end{proof}

\begin{remark}[What is paid, and what is not being claimed]
The factor $e^{t_0-t_m}=A_q^{-1}q^{c_m-10}e^{o(1)}$
explicitly pays the possible strong-rate growth on this interval.
It is not hidden in an order-one constant.
Before $t_m$, all-label early comparison and the exact target
bound pay the full radial evolution.
We therefore do not need an $O(1)$ estimate for Q3SB's
\texttt{unpaidradial} integral, and do not claim to have proved that
sharper integral estimate. In particular we do not extend Q2EC's
off-chart transport to Q3 neighbours by assumption.
There is no enlargement of $G_+$ or change to the partition used by
LCRC, S2CP or CFA. The price of this simpler bound is the weaker
positive exponent in Theorem~\ref{an03pisbv1:main}.
\end{remark}

\section{The regularized post-entry gain and the endpoint check}
\begin{proposition}[Absolute gain from $t_0$]
\label{an03pisbv1:postgain}
For every original receiver define
$N(t)=|\sin\theta(t)|+|\sin\psi(t)|$, $N_0=N(t_0)$.
There is a fixed $A_N\ge1$ such that
\begin{equation}
 \frac{m(t,b)}{m(t_0,b)}
 \le e^{2B+\eta_q+1/\beta}
   \max\{1,A_N(N_0+q)^2e^{\beta I(t)}\}^{1/\beta},
 \quad
 \eta_q=C(1+M_{\rm res})qL_T\ell .
 \label{an03pisbv1:regularizedgain}
\end{equation}
Here $1/\beta=1+q^2/\lambda$ exactly. It is at most $1+\nu$
for the fixed admissible $\nu$ once $q$ is small.
The simpler simultaneous bound is
\begin{equation}
 m(t,b)\le m(t_0,b)e^{2B+\eta_q}e^{I(t)} .
 \label{an03pisbv1:capone}
\end{equation}
The latter bound has no $A_N$ factor.
\end{proposition}
\begin{proof}
GESG's global matrix-value sine inequality gives
\[
 D^+N\le\{r+X+I'/2+a_*q\}N+a_*q,\qquad
 a_*=C(1+M_{\rm res}).
\]
Replace $I'/2$ by $\beta I'/2$ at an integrated cost
$O(q^2L_T\ell)$, and note
$\beta I'/2\ge a_0:=\lambda_{\min}(1-c_b)/2>0$.
Duhamel, \eqref{an03pisbv1:gate} and $\int r<B$ give
\[
 N(t)\le\min\{2,C_N(N_0+q)e^{\beta I(t)/2}\},\qquad
 C_N=e^{B+G+1}\max(1,a_*/a_0)
\]
for sufficiently small $q$, after harmless enlargement.
Take $A_N=\max(1,C_N^2/4)$.

DECG's signed weak-gate estimate, including its negative-input
case, gives directly from \eqref{an03pisbv1:exact}
\[
 (\log m)'\le2r+C(1+M_{\rm res})q
                     +I'(\sin\theta\sin\psi)_+ .
\]
Since $(\sin\theta\sin\psi)_+\le\min(1,N^2/4)$, TAILREP's exact
gain-integration lemma proves \eqref{an03pisbv1:regularizedgain}.
Retaining only the cap one proves \eqref{an03pisbv1:capone}.
This also explains why no additional gate-tail sign assumption
or chart entry is needed after $t_0$.
\end{proof}

\begin{remark}[The endpoint $b=\pi$]
The endpoint is included in Lemmas~\ref{an03pisbv1:domination}--
\ref{an03pisbv1:bridge}, with equality in the target comparison.
Its nonzero Gaussian seed and crossing time $4\ell/\lambda+O(1)$
are kept explicitly. Q3SB's lower bounds
\[
 \frac{m(t_m,\pi)}{q^{10}}
       \ge c e^{-Cq^u}q^{2-\lambda c_m},\qquad
 \frac{m(t_0,\pi)}{q^{10}}
       \ge c A_q^{-\lambda}q^{2-10\lambda}
\]
are consistent with the new upper bounds.
At $t_m$ the first lower bound is represented by the
$q^{12-\lambda c_m}$ term.
At $t_0$,
\[
 (12-10\lambda)-[2+(1-\lambda)c_m]
      =\frac{11-19\lambda}{2}>0
\]
on the box, so the coarser entry bound is indeed larger.
In the post-entry bound the additive $q$ remains even if $N_0=0$.
Neither a zero initial sine nor exact original-flow angular ordering
is substituted for the created seed.
\end{remark}

\section{Closing the \texorpdfstring{$\pi$}{pi} mass budget}
\begin{proof}[Proof of Theorem~\ref{an03pisbv1:main}]
Only the inner zone remains after \eqref{an03pisbv1:frozenbudget}.
For each of its labels, combine
\eqref{an03pisbv1:entrymass},
\eqref{an03pisbv1:capone} and the stopped clock
\eqref{an03pisbv1:clock}. At every endpoint,
\[
 m(t,b)\le
 C e^{2B+\Xi_q}E_gA_q^{\lambda-1}
       \left(q^{c_m-k}
             +q^{2+(1-\lambda)c_m-k}
             +q^{4/\lambda-2-k}\right).
\]
The powers are exactly the three $\delta_j$ in
\eqref{an03pisbv1:powers}, since
$(1-\lambda)c_m=4+(1-\lambda)/2$.
The initial width of this zone is $qs_q$, and its measure is
$qs_q/(2\pi)$. This proves \eqref{an03pisbv1:pibudget}.

For completeness the rational margins need no relation beyond
the ratio box. $\delta_1$ increases with $\lambda$, and
\[
 \delta_1\ge\frac{15259}{9240}>\frac{411}{800}.
\]
The increasing affine function $\delta_2$ has minimum
$411/800$ at $\lambda=169/400$.
Finally $\delta_3'=10-4/\lambda^2<0$ on the box, so
\[
 \delta_3\ge\frac{521}{490}>\frac{411}{800}.
\]
All factors outside these powers are fixed or tend to one,
$A_q$ stays bounded above and below, and
$s_q=O(\sqrt{\log\ell})$. Multiplication by the duration
$L_T\ell$ still tends to zero. This closes the supremum and
integral assertions.
\end{proof}

\section{The remaining \texorpdfstring{$3\pi/2$}{3pi/2} consumer and the constant order}
Only the fixed initial inner ancestry
\[
 J_3=\{3\pi/2-qs_q<b\le3\pi/2\}
\]
remains unpaid on the other side. No claim about its entry, its
radial lag or its relative shape is proved here.

For an already charted part, S2CP's \texttt{latestrip} consumer is a
bounded actual-reference entry chart and the \emph{raw} law
\[
 \eta_{\rm late}\{\xi>v\}
     \le C_T\min\{1,(\epsilon_h/v)^p\},\qquad
 \dd\eta_{\rm late}=m(t_0,b)\dd\lambda_0,\qquad p=3/s_{\rm pr},
\]
with the corresponding TAILREP cutoff/moment. Here
\[
 d=(1-\lambda P_\rho)/2,\quad s_{\rm pr}=d/\omega,\quad
 \alpha=\lambda P_\rho/2,\quad
 \epsilon_h=\sqrt2 A_q^d(1-M_0)^\alpha q^{10d-s_{\rm pr}} .
\]
These are the source's profile parameters, not the offsets $s_0$.

If entry is later, let $\tau_b\ge t_0$ be the actual joint chart
arrival, with $\tau_b=t_0$ for labels already charted, and
$\tau_b=\infty$ otherwise. Let $\xi_b$ be its scaled displacement
from the actual reference at arrival. A sufficient exact consumer
is an arrival-valid gain dominating current mass by
\begin{align}
 \mathcal W_3(t)={}&
 \int_{J_3\cap\{\tau_b>t\}}m(t,b)\dd\lambda_0(b)\nonumber\\
 &+C_m\int_{J_3\cap\{\tau_b\le t\}}m(\tau_b,b)
 \max\{1,A_0q^2\xi_b^2e^{\beta[I(t)-I(\tau_b)]}\}^{1+\nu}
                      \dd\lambda_0(b),             \label{an03pisbv1:lateconsumer}
\end{align}
together with
\[
 \sup_{[t_0,t_*]}\mathcal W_3/q=o(1),\qquad
 \int_{t_0}^{t_*}\mathcal W_3/q=o(1),
\]
or the second inequality with CFA's small fixed allowance.
Both the arrival-valid gain and this weighted estimate remain to
be proved for $J_3$. They are not assumed in the $\pi$ theorem.
The unarrived term must not be dropped.

The new constants have the following dependencies.
Loss/Gaussian constants, $k_G$, $C_{\rm pre}$, the target root
constants, the crossing constants, $c_A,C_A$, and the fixed
mass-bridge cost $e^6$ are model/S1 constants, independent of
LCRC guards. $C_{\rm dur}$ uses the retained duration.
The middle-freezing constants depend on $\varepsilon$,
$C_{\rm dur}$ and $1+C_{\rm pre}+M_{\rm res}$, all fixed before $q$.
For the optional sharper post-entry inequality,
$C_N,A_N$ depend on $B,G,M_{\rm res},c_b$ as displayed.
They do not enter the main estimate through any new $q$ exponent.
The main $\pi$ estimate has the explicit fixed factor $e^{2B}E_g$,
the power $A_q^{\lambda-1}$ and the vanishing clock/rate defect
$\Xi_q$; its remaining constants are model/S1 constants.
No chart radius, core mass bound or entry threshold is enlarged.

Thus the existing S2CP/CFA constant order is unchanged:
model and guard-independent core constants, fixed stopped guards
and field allowances, $C_0$ and its entry constants as in those
sources, the auxiliary $\varepsilon$, and finally $q$.
For any already fixed CFA allowance, the proved $\pi$ contribution
and both frozen-zone contributions are smaller than that allowance
for sufficiently small $q$. The $3\pi/2$ inner contribution still
has to be paid before CFA's \emph{whole-strip} input is available.
We stop here: no joint weak-passage bootstrap, no CLOCK passage
application, no new entry contracts, and no consolidation.
\end{document}

